\section{LoRA in Continual Learning: O-LoRA}
\subsection{Introduction}
\subsubsection{Continual learning}
Continual learning (CL) lets a model learn new tasks sequentially while preserving knowledge of old tasks. This paradigm aims to give AI systems some of the plasticity of biological systems, allowing them to accumulate and use knowledge in changing environments. Continual learning is important for AI systems that must adapt to new tasks without repeatedly retraining on all previous data.

Three common CL settings are:
\begin{itemize}
	\item \textbf{Task-incremental learning (TIL)}: the model learns a sequence of tasks, with the task ID known at both training and inference time. For example, it might be given sentiment analysis for the first task, then review rating for the next.
	\item \textbf{Class-incremental learning (CIL)}: the model learns new classes sequentially, with the task ID known during training but not given at inference time. The model must predict among all classes seen so far, which makes this a harder setting than TIL.
	\item \textbf{Domain-incremental learning (DIL)}: the model learns the same task under changing data distributions, without needing a task ID at inference time. For example, it might learn sentiment analysis first on music reviews, then on hotel reviews.
\end{itemize}

\subsubsection{Problem}
One fundamental challenge of continual learning is catastrophic forgetting: after learning new tasks, a model may lose performance on old ones. This creates a trade-off between stability (retaining old knowledge) and plasticity (learning new knowledge). Various methods have been devised to address this issue. Regularization-based methods add terms to the training loss that constrain new updates, with the aim of preserving old task knowledge. We will look at one such method called O-LoRA.

\subsubsection{Solution}
O-LoRA is a regularization-based CL method built on the idea that keeping task update subspaces orthogonal can reduce interference with old tasks. It gives each task a separate LoRA update, freezes the updates of previous tasks, and adds a penalty that encourages the subspace of new task \(t\) to be orthogonal to those of previous tasks \(i<t\). These notes follow \href{https://arxiv.org/abs/2310.14152}{Wang et al., \emph{Orthogonal Subspace Learning for Language Model Continual Learning}}.

\subsection{Formulation}

\subsubsection{Task-specific LoRA updates}

Suppose tasks \(\mathcal{D}_1,\ldots,\mathcal{D}_T\) arrive sequentially, and only the current task's data \(\mathcal{D}_t\) are available when learning task \(t\). For a pre-trained weight matrix \(W_0 \in \mathbb{R}^{d \times k}\), O-LoRA learns a separate low-rank update \(A_t B_t\) for each task, where \(A_t \in \mathbb{R}^{d \times r}\), \(B_t \in \mathbb{R}^{r\times k}\), and \(r \ll \min(d,k)\). After task \(t\), the adapted matrix is
\[
	W_t = W_0 + \sum_{i=1}^{t} \Delta W_i
\]
where \(\Delta W_i = \frac{\alpha}{r} A_i B_i\).
With input \(x\in\mathbb{R}^{k}\), the adapted layer computes
\[
	h_t = W_t x
\]
The base matrix and the previous task factors \(\{A_i,B_i\}_{i<t}\) remain fixed while \(A_t\) and \(B_t\) are trained.
The complete model update \(\Delta\Phi_i\) collects these updates across all selected weight matrices. Thus,
\[
	\Phi_t = \Phi_0+\sum_{i=1}^{t}\Delta\Phi_i
\]
Here, \(\Phi_t\) denotes the complete model parameters after incorporating the updates from tasks \(1,\ldots,t\), while \(W_t\) denotes an individual adapted weight matrix.
For brevity, the matrix-level formulation is shown for one selected weight matrix; the same updates and orthogonality penalties are applied to all selected matrices.
Then the task loss is
\[
	\mathcal{L}_{\mathrm{task}} =
	- \sum_{(x,y)\in\mathcal{D}_t}\log P_{\Phi_t}(y\mid x)
\]

\subsubsection{Orthogonality constraint}

O-LoRA approximates the task-specific update subspace \(\mathcal U_t \subseteq \mathbb{R}^{d}\) by the column space of \(A_t\) for task \(t\). To make \(\mathcal{U}_t\) orthogonal to a previous task's subspace \(\mathcal{U}_i\), O-LoRA would require \(A_i^{\top}A_t=0\). It encourages this condition with the penalty
\[
	\mathcal{L}_{\mathrm{orth}}(A_t)
	= \sum_{i=1}^{t-1} \left\|A_i^{\top} A_t\right\|_F^2
	= \sum_{i=1}^{t-1} \sum_{p=1}^{r}\sum_{q=1}^{r} \left((A_i^{\top} A_t)_{pq}\right)^2
\]
where \(\|\cdot\|_F\) denotes the Frobenius norm of a matrix.

For a positive weight \(\lambda\), the current task's training objective can be written as
\[
	\min_{A_t,B_t}
	\mathcal{L}_{\mathrm{task}}
	+ \lambda \mathcal{L}_{\mathrm{orth}}(A_t)
\]
The model fits the current task while penalizing overlap with previous update subspaces. A finite penalty encourages orthogonality but does not guarantee it.\footnote{The paper writes positive log-likelihood alongside a positive orthogonality term; the negative log-likelihood form here makes the intended penalty explicit.}

\subsubsection{Training and inference}

In the paper, LoRA is applied to the attention query and value weights. Training task \(t\) needs the current data and the frozen earlier LoRA factors, but no stored examples from previous tasks. The number of task factors grows during training. For inference, their updates can be merged into the pre-trained weights as in \(W_t\), so the model can make predictions without selecting an adapter using a task ID.
